\section{From an independent set to bounded subchain lists}
\label{sec:lists}

Fix $\sigma\in(0,1)$ and assume that the Label Cover value is at most $\sigma$.
The clique branch already has the required soundness, so suppose
it does not occur. Fix a nonempty independent vertex set $\mathcal A$
in the output graph. We first turn it into functions on the tuple tori,
then approximate their restrictions to small subchains using bounded
coordinate lists. Source soundness will show that lists belonging to
separated subchains are usually disjoint. Throughout the analysis, set
$L=64$.

\subsection{Odd functions and compatibility}

Let $A\subseteq\mathscr X$ be the set of locations of vertices in
$\mathcal A$, with multiplicities removed. Then
\begin{equation}\label{eq:independent-separation}
 \dist(A,TA)>1/8.
\end{equation}
For locations of distinct selected vertices this follows from the missing
edge in \eqref{eq:edge}. For a point compared with itself, it follows
from exclusion of \eqref{eq:clique-test}. Since the sets are finite,
these pointwise strict inequalities give \eqref{eq:independent-separation},
including when distances are infinite.

On $\mathscr X$ define
\begin{equation}\label{eq:bumps}
 f(x)=\bigl(1-32\dist(x,A)\bigr)_+
       -\bigl(1-32\dist(Tx,A)\bigr)_+,
 \qquad (u)_+=\max\{u,0\},
\end{equation}
with each bump defined to be zero at infinite distance.
This function takes values in $[-1,1]$, satisfies $f(Tx)=-f(x)$,
and equals $1$ on $A$.
It is constant across every zero link.
Within each grid copy it is $64$-Lipschitz for the sup circle metric:
distance to $A$ is $1$-Lipschitz on a component meeting $A$, and its
bump is identically zero on a component not meeting $A$.
Also the shortest-path pseudometric never exceeds a within-copy
link distance.

At each tuple $\mathbf q$, extend its grid function to the whole
torus $\T^{M_i}$. Using the real-valued Lipschitz extension formula of McShane
\cite{McShane1934}, set
\[
 g_{\mathbf q}(x)=
 \min_{y\in\Gamma^{M_i}}
       \bigl(f(y)+L\,d_\infty(x,y)\bigr),
\]
where $y$ ranges over that tuple's copy and $d_\infty$ is the sup circle
metric. This is $L$-Lipschitz and equals $f$ on the grid.
Clip its values to $[-1,1]$, obtaining $\widetilde g_{\mathbf q}$, and put
\[
 F_{\mathbf q}(x)=
 \frac{\widetilde g_{\mathbf q}(x)
       -\widetilde g_{\mathbf q}(Tx)}2.
\]
The resulting function is still an extension of the odd grid data,
is $L$-Lipschitz, is valued in $[-1,1]$, and is odd under $T$.
Fix this entire family of functions before sampling any test chain.
Along a chain write $F_i=F_{\mathbf q_i}$.

For every chain and every $i\le j$, these extensions satisfy
\begin{equation}\label{eq:compatibility}
 |F_j(x)-F_i(x\circ\pi_{ij})|
 \le 2L/D,\qquad x\in\T^{M_j}.
\end{equation}
Indeed, round each coordinate of $x$ once to a grid point $y$
within $1/D$. Pullback does not enlarge this error.
The grid points $y$ and $y\circ\pi_{ij}$ have equal $f$ values by
the chain of zero links, so the two extension errors total at most
$2L/D$. Intermediate layers contribute no additional approximation
error.

\subsection{Bounded lists for a subchain}

The functions $F_i$ now supply the analytic data we need along any chain.
For each small set of layers, we will choose a bounded list of answer
coordinates that approximates every allowed coefficient choice.
The bound must be independent of the Label Cover alphabets, and the
choice must depend only on the data of that subchain.

Fix an approximation tolerance $\gamma>0$.
For $\varnothing\ne I\subseteq[r]$ with $|I|\le s$, put $b_I=\min I$.
Given $t\in\mathcal T_m^I$, define
\begin{equation}\label{eq:subchain-functions}
 h_{I,t}(\theta)
 =F_{b_I}\left(
  \left(\sum_{j\in I}t_j\theta_{j,\pi_{b_Ij}(\kappa)}
            \pmod1\right)_{\kappa\in M_{b_I}}\right),
 \qquad
 \theta\in\prod_{j\in I}[0,1]^{M_j}.
\end{equation}
The entries of $\theta$ are independent uniform real variables.
This evaluates $F_{b_I}$ at the continuous version of the sampled
location in \eqref{eq:location}, using only the layers of $I$.
The phase map has Lipschitz constant at most $|I|\le s$:
coordinate duplication by a projection cannot enlarge a sup distance.
Thus $h_{I,t}$ is $sL$-Lipschitz, independently of the label-set sizes
and projection fibers.

\begin{lemma}[Bounded subchain lists]\label{lem:subchain-lists}
For every fixed occurrence chain and every nonempty $I\subseteq[r]$
with $|I|\le s$, there are sets
$S_I^j\subseteq M_j$, $j\in I$, with
\begin{equation}\label{eq:list-size}
 \sum_{j\in I}|S_I^j|\le d,
 \qquad d:=\max\{1,J(sL,\gamma)(m+1)^s\},
\end{equation}
such that, for every $t\in\mathcal T_m^I$, a junta $g_{I,t}$ using
only scalar coordinates $(j,\ell)$ with $\ell\in S_I^j$ satisfies
\begin{equation}\label{eq:subchain-junta}
 \|h_{I,t}-g_{I,t}\|_2<\gamma.
\end{equation}
The lists can be chosen as functions only of the subchain's label sets,
internal projection maps, and functions $F_j$, $j\in I$.
Identical subchain data receive identical choices.
\end{lemma}

\begin{proof}
For each of the at most $(m+1)^s$ coefficient choices, apply
Theorem~\ref{thm:junta} to $h_{I,t}$ with error $\gamma$.
This uses at most $J(sL,\gamma)$ scalar coordinates.
For each block $j$, take the union of its selected coordinates over
all coefficient choices to obtain $S_I^j$ and \eqref{eq:list-size}.

To make the choice depend only on the stated subchain data, order the
scalar coordinates and choose the first subset of size at most
$J(sL,\gamma)$ whose coordinate conditional expectation has error
less than $\gamma$. The theorem guarantees such a subset.
These choices enter only the soundness proof and are never computed
by the reduction. In particular, a list includes approximating
coordinates for every coefficient vector before any coefficients
are sampled.
\end{proof}

\subsection{Why the lists can be decoded locally}

Use the fixed rule of Lemma~\ref{lem:subchain-lists} along a uniform
occurrence chain.
Using $M_r=\Sigma_R^{r-1}$ as a common label universe, set
\begin{equation}\label{eq:terminal-lists}
 A_I=\bigcup_{j\in I}\pi_{jr}(S_I^j),
 \qquad \varnothing\ne I\subseteq[r],\quad |I|\le s.
\end{equation}
By \eqref{eq:list-size}, each $A_I$ has size at most $d$.
We first control intersections between lists on separated sets of layers.

\begin{lemma}[Separated lists]\label{lem:separated-lists}
For every fixed $I,J$ with $1\le|I|,|J|\le s$ and
$\max I<\min J$,
\[
 \Pr_{\text{chain}}(A_I\cap A_J\ne\varnothing)\le d^2\sigma.
\]
\end{lemma}

\begin{proof}
Put $h=\max I<\min J$ and condition on every occurrence except $c_h$.
The remaining occurrence is still uniform in the original instance.
At each layer in $I$, position $h$ of the tuple is the left question
$u_h$. Every internal projection between layers of $I$ uses only
tests with index less than $h$. Thus all subchain data for $I$,
including its already fixed tuple functions, depend only on $u_h$
and the conditioned background, not on the opposite question or on
the projection of $c_h$.
Likewise, all data for $J$ depend only on $v_h$ and the background.
Figure~\ref{fig:locality} illustrates this separation.

\begin{figure}[t]
\centering
\begin{tikzpicture}[>=Stealth, every node/.style={font=\small}]
\node (a) at (0,0) {$(u_1,u_2,u_3)$};
\node (b) at (3.8,0) {$(v_1,u_2,u_3)$};
\node (c) at (7.6,0) {$(v_1,v_2,u_3)$};
\node (d) at (11.4,0) {$(v_1,v_2,v_3)$};
\draw[->] (a) -- node[above] {$c_1$} (b);
\draw[->] (b) -- node[above] {$c_2$} (c);
\draw[->] (c) -- node[above] {$c_3$} (d);
\draw[dashed,gray] (5.7,-0.9) -- (5.7,1.1);
\node at (1.9,-0.65) {Position $2$: own question $u_2$};
\node at (9.5,-0.65) {Position $2$: own question $v_2$};
\node[font=\footnotesize] at (1.9,0.9) {left subchain};
\node[font=\footnotesize] at (9.5,0.9) {right subchain};
\end{tikzpicture}
\caption{Four layers with the separating test $c_2$ left unconditioned.
Once $c_1,c_3$ are fixed, the left and right subchain data depend only
on the respective own question of $c_2$. Answer maps act in the position
indicated by each test.}
\label{fig:locality}
\end{figure}

Form a list of candidate left answers by taking position $h$
of every label in every $S_I^j$ for $j\in I$.
This gives at most $d$ elements of $\Sigma_L$, as a function of $u_h$
and fixed background. The analogous list for $J$ gives at most $d$
elements of $\Sigma_R$ as a function of $v_h$.
Choose deterministic orderings and pad each list to length $d$
with an arbitrary answer, also when the list is empty.

If $A_I$ and $A_J$ intersect, some left and right selected labels have
the same image at layer $r$. Equality in position $h$ says that
$\pi_{c_h}$ sends the associated left candidate to the right candidate:
on the left, exactly the test $c_h$ acts in that position, and on the
right the position is already a right-answer coordinate.
For each pair of list positions, selecting those answers for each
own question is a deterministic strategy for the original Label Cover
instance. Its success probability over the uniform $c_h$ is at most
$\sigma$. A union bound over the $d^2$ position pairs proves the
conditional bound, and averaging over the background proves the lemma.
\end{proof}

The fixed functions may depend on the entire instance and on
$\mathcal A$. This causes no loss of locality: they are fixed before
the sampling experiment, so evaluating the family at a tuple reveals
only that tuple. In particular, if two occurrences have the same own
question but different opposite questions or projections, the list
rule at that side uses identical data.

We have obtained bounded lists in a common answer set, and source
soundness controls intersections of lists on separated layer sets.
The next section turns this separation into a stronger property:
each listed answer can be assigned to one layer common to all of its
occurrences inside the sampled set $B$.
